\documentclass[10pt]{article}

% Specify the margins and text size.
\setlength{\textwidth}{6.4in}
\setlength{\textheight}{9.5in}
\setlength{\oddsidemargin}{0pt}
\setlength{\evensidemargin}{0pt}
\setlength{\topmargin}{0pt}
\setlength{\hoffset}{.05in}
\setlength{\voffset}{-1in}

\setlength{\parskip}{5pt}
\setlength{\parindent}{0pt}

% Load some fonts and symbol packages
\usepackage{latexsym}
\usepackage{pifont}       % contains 'star' symbol for counterinsurgency handbook title
\usepackage{yfonts} 
\usepackage{amsmath}
\usepackage{amsfonts}

\usepackage{graphicx}     % actually, this is loaded by pstricks
\usepackage[T1]{fontenc}
\usepackage{ifthen}
\usepackage{pstricks,pst-grad,pst-text,pst-node,multido,pst-plot,calc,pst-3dplot}
%\usepackage[all]{xy}
%\usepackage{animate}

% The hyperref package inserts links into the pdf.
\definecolor{MyLinkColor}{rgb}{.1,.2,1}
\definecolor{MyCiteColor}{rgb}{.1,1,.2}
\definecolor{MyURLColor}{rgb}{.4,.4,.4}
\usepackage[backref=true,pagebackref=false,hyperindex,colorlinks=true,
  linkcolor=MyLinkColor,urlcolor=MyURLColor]{hyperref}


% The tweaklist package is something I found on the web.  It provides a simple interface
% for making changes to spacing used in the itemize and enumerate environments.  Comment
% this out if you don't care to use tweaklist.
\usepackage{tweaklist}
\renewcommand{\itemhook}{\setlength{\parskip}{2pt}\setlength{\parsep}%
{1pt}\setlength{\topsep}{0pt}\setlength{\itemsep}{0pt}}

\newcommand{\U}{\underline{\hspace{5pt}}}

\usepackage{listings}
\newcommand{\Z}{\hphantom{0}}

\begin{document}

\pagestyle{empty}

\href{http://www.shepherd.edu}{\includegraphics[height=1.75cm]{logo-high-res.eps}}
\vspace{-1.69cm}

{\small
\begin{tabular}{cl}
& Math 314\\
& Statistics\\
\hspace{5.28in} & %19 October 2011
\end{tabular}
}

\begin{center}
\textbf{\large  Exam II}
\end{center}
No books or notes  may be used on this exam.
\medskip

1. (30 points) In a certain class, midterm scores average out to 70 with an SD of 15, as do 
scores on the final.  The correlation between midterm scores and final scores is 
about~$0.50$.   

\hspace{20pt} a) Write down the equation for the regression line. %(10 points)
\vspace{1.0in}

\hspace{20pt} b) Sketch the regression line. %(5 points)
\vspace{1.25in}

\hspace{20pt} c) 
Estimate the average final exam scores for students whose midterm scores were
85, 40 and 70. %(10 points)
\vspace{2.75in}

\hspace{20pt} d) Estimate the midterm score for a student who scored 90 on the final exam.
%(5 points)
\vfill
\eject
{\ }


%\#7 page 144 (word problem about correlation)
%\#8 page 155 (word problem about correlation)
2. (10 points) 

\hspace{20pt} a) For a representative sample of cars, would the correlation between the
age of the car and its gasoline economy (miles per gallon) be positive or negative?  Explain
\vspace{1.5in}

\hspace{20pt} b ) The correlation between gasoline economy and income of owner turns out to 
be positive?  How do you account for this association?
\vspace{1.5in}

3. (10 points) A box contains 5 red marbles and 4 green ones.   

\hspace{20pt} a) If three marbles are drawn \textit{with replacement}, what is the chance
that two of the three are red?
\vspace{1.2in}

\hspace{20pt} b) If three marbles are drawn \textit{without replacement}, what is the
chance that two of the three are red?
\vspace{1.2in}

4. (10 points) Cards are drawn at random without replacement from a 
well-shuffled deck. Calculate the following probabilities.
\medskip

\hspace{20pt} a) First card is a $\clubsuit$ or a $\spadesuit$
\bigskip\bigskip

\hspace{20pt} b) Tenth card is a $\heartsuit$ or a face card
\bigskip\bigskip

\hspace{20pt} c) Second card is an ace given that the first one drawn was a face card
%\vspace{1in}



%4. (10 points) Test scores on a certain exam average out to 75 with a standard 
%deviation of 3.  One student scored 84.   Was that a very high score?  Moderately high?  
%Or about average?  Justify your answer.
\vfill
\eject
{\ }
5. (10 points) For people age 25 and over in the U.S. in 2005, the relationship between
age and educational level (years of schooling completed) can be summarized as
follows. 
\begin{align*}
\mbox{average age}         & \approx \mbox{50 years,}    & \mbox{SD} & \approx \mbox{16 years}\\[3pt]
\mbox{average education level}  & \approx \mbox{13.2 years,}  & \mbox{SD} & \approx \mbox{3.0 years,
   \hspace{15pt}} r\approx -0.2
\end{align*}

\hspace{20pt} a) True or false:  For people age 25 and older in the U.S. in 2005, older people tended 
to be less\vspace{-5pt}

\hspace{20pt} \hphantom{a) }  educated than younger ones.
\bigskip

\hspace{20pt} b) True or false, and explain:  as you get older, you become less educated.  
If this statement is false,\vspace{-4pt}

\hspace{20pt} \hphantom{a) }
what could account for the negative correlation?  
\vspace{2in}


6. (15 points) Suppose that the regression line of $y$ on $x$ 
is \[y=\left(\frac{\sqrt{3}}{2}\right)\,\left(\frac{8}{\sqrt{3}}\right)\,(x-1)\] 
where $\mbox{mean}_x=1$, $\mbox{mean}_y=0$, $\mbox{SD}_x=\sqrt{3}$, $\mbox{SD}_y=8$
and $r=\sqrt{3}/2$.\vspace{3pt}

\hspace{20pt} a) Calculate the RMS error for regression 
$\mbox{RMS}=\mbox{SD}_y\,\sqrt{1-r^2}$.  What can you say about the RMS\vspace{-4pt}

\hspace{20pt}\hphantom{b) } for the SD line?
\vspace{2in}

\hspace{20pt} b) If $x=4$, then we expect 68\% of the $y$-values to be
between \ \underline{\hspace{60pt}} \ \ and \  \underline{\hspace{60pt}} \ .
\vspace{2in}



\vfill
\eject

(15 points)

7. Consider the following data.
\begin{tabular}{|c|c|}\hline
$x$ & $y$\\\hline
0   & 2\\
1   & 4\\
2   & 3\\\hline
\end{tabular}

\hspace{10pt} a) Make a scatter plot of the data.  
\vspace{3.5in}

\hspace{10pt} b) An investigator uses the line $y=2\,x + 1$ to predict $y$ from $x$.
Add the line to your sketch.
\medskip

\hspace{10pt} c) Calculate the RMS error for the line:
$\displaystyle\sqrt{\frac{1}{n}\sum\left(\mbox{predicted $y$} - \mbox{actual $y$}\right)^2}$.
\vspace{2.8in}
\vfill
\eject

\end{document}
